\section{Background}
\label{sec:background}

The Bailout Protocol addresses a structural gap in the governance of autonomous multi-agent AI systems: when a spawned agent encounters a condition it cannot safely resolve, existing frameworks provide no architecturally defined mechanism for routing that exception to a human decision-maker without relying on intermediate machine actors that may themselves be compromised, misconfigured, or insufficiently authoritative. This section situates the Bailout Protocol within four intersecting research threads — accountability in recursive delegation chains, human-in-the-loop design principles, fail-safe architecture in safety-critical systems, and the BX3 Framework's three-layer model as an architectural substrate for structured exception escalation.

% -------------------------------------------------------------------------- %
\subsection{Accountability in Multi-Agent AI Systems}
\label{sec:accountability}

Modern AI agent frameworks permit recursive delegation: a parent agent spawns a child to execute a sub-task, which may in turn spawn a grandchild, forming a delegation chain of arbitrary depth. This architectural pattern is powerful because it enables scalable, asynchronous task decomposition, but it creates a corresponding accountability deficit. As the chain lengthens, the causal and legal origin of any given action becomes progressively more difficult to establish. A human principal authorises the top of the chain, but intermediate machine actors are interposed between that principal and the eventual action — and those intermediate actors are not legal persons capable of bearing liability.

This problem is not hypothetical. Hood~\cite{hood2025} identifies \emph{authority laundering} as a core failure mode in deep delegation chains: when a chain is sufficiently long, the origin of authorisation becomes obscured. An agent at depth $n$ takes an action that appears to originate from depth $n-1$, which traces back through depth $n-2$, and so on, eventually reaching what is claimed to be an authorised human principal. If any link in this chain is mutable — if a parent pointer can be redirected after an action occurs — then the chain of authority is retrospectively corruptible. The action at depth $n$ could be re-parented to a different ancestor after the fact, creating a plausible but false narrative of authorised delegation. This is not a policy failure; it is a structural failure of the chain itself.

The problem is compounded by the probabilistic nature of AI inference. Because the same input may produce different outputs on different invocations, a delegation chain that was trustworthy at provisioning may become unreliable as child agents drift — through model updates, context corruption, or adversarial manipulation — into behaviours not anticipated by the parent. If the parent pointer is mutable, this drift is invisible. If the chain is immutable, the drift is attributable: the original parent remains on record as responsible regardless of the child's current state.

Existing accountability infrastructure — audit logs, access control policies, approval workflows — assumes a single authoritative principal making decisions in a centralised context. These assumptions break down in recursive multi-agent systems, where the chain of principals may be deep, the actions may be asynchronous, and the human at the root may have no direct knowledge of what any given child agent is doing. Audit logs are only as reliable as the system that writes them; if the writing process itself can be manipulated by a compromised intermediate actor, the audit trail records the attack rather than the original authorised intent.

The accountability gap is therefore not merely a governance concern — it is an architectural one. Policy-level accountability measures cannot close a gap that is created by the structural properties of the delegation chain itself. Immutable parent chains, deterministic execution logging, and architecturally enforced exception escalation paths are required to make accountability tractable in deep recursive agent systems.

% -------------------------------------------------------------------------- %
\subsection{Human-in-the-Loop Design}
\label{sec:hitl}

The human-in-the-loop (HITL) paradigm is the primary mechanism by which human oversight is maintained over autonomous systems. The core principle is straightforward: a human principal retains the ability to approve, override, or abort any action proposed or taken by a machine actor. In practice, however, HITL implementation spans a wide spectrum of sophistication, and the specific design choices made at this level determine whether human oversight is structurally effective or merely nominal.

At the most basic level, HITL appears as a watchdog: a human monitors the outputs of an autonomous system and intervenes when those outputs appear incorrect or unsafe. This model is common in contemporary AI assistant deployments, where outputs are presented to a user for review before any consequential action is taken. However, watchdog HITL scales poorly. A human monitoring a single agent's outputs can maintain adequate oversight; a human monitoring a delegation chain of arbitrary depth, with children spawning children asynchronously, cannot. The human becomes a bottleneck, and the operational incentive to reduce that bottleneck creates pressure to reduce the fidelity of human review.

More structured HITL models introduce tiered approval thresholds. Low-stakes actions execute autonomously; medium-stakes actions require approval from a designated human reviewer; high-stakes actions require approval from an authorised principal with legal standing. This model is familiar from financial systems and nuclear plant operations, where authorisation matrices define who may approve what class of action. However, tiered approval still depends on the machine actor correctly identifying the stakes of its own proposed action — a classification problem that is itself susceptible to misalignment and adversarial manipulation.

Dijkstra observed that the quality of a system is constrained by the quality of its control structures~\cite{dijkstra1982}. In the context of HITL, this implies that the effectiveness of human oversight is constrained by the quality of the interface between machine actors and human principals. If that interface is a flat notification channel, human oversight degrades under load. If it is a structured escalation hierarchy with well-defined trigger conditions and bounded resolution time, human oversight can be engineered to remain reliable under operational stress.

The design of the escalation path is as important as the existence of one. An escalation mechanism that routes through intermediate machine actors before reaching a human is not a HITL mechanism — it is a multi-tier machine delegation system with a human at the end of a chain. A true HITL mechanism excludes intermediate machine actors from the critical path when human decision-making is required. The Bailout Protocol is designed around this principle: its escalation chain is defined to terminate at the first Human Accountability Anchor, bypassing all machine intermediaries regardless of their position in the delegation hierarchy.

% -------------------------------------------------------------------------- %
\subsection{Fail-Safe Design in Safety-Critical Systems}
\label{sec:failsafe}

Safety-critical systems — aviation, nuclear power, surgical robotics, autonomous vehicles — have long confronted the challenge of designing autonomous agents that operate reliably in environments where failure has catastrophic consequences. Thefail-safe design tradition in these domains offers important lessons for multi-agent AI governance.

A fail-safe system is one that fails in a manner that preserves human safety, rather than one that simply avoids failure. In aviation, fly-by-wire systems are designed to enter a controlled glide and surrender control to the pilot if a flight computer fails — they do not continue operating on the basis of stale or corrupted sensor data. The underlying principle is that the system must have a defined, safe state to which it reverts when its operating assumptions are violated. This is not merely a reliability goal; it is an architectural requirement that must be satisfied before the system is deployed.

Levesitt~\cite{tristr81} introduced the concept of transactions as an abstract unit of correct computation — a sequence of operations that transforms the system from one consistent state to another, with atomic commit or rollback semantics. This principle directly applies to multi-agent systems: an agent's action should be treated as a transaction whose outcome is either fully committed or fully rolled back, with no intermediate state in which the system's consistency is compromised. Fail-safe rollback is the architectural expression of this principle in the context of autonomous operation.

In autonomous vehicles, Bansal and segment~\cite{bansal2019} demonstrated that even well-designed perception systems can fail catastrophically in adversarial conditions — specifically, they showed that adversarial physical perturbations to road signs can cause state-of-the-art classifiers to misread speed limits with high confidence and at long range. The practical implication is that no perception system, however well-trained, can be trusted to correctly classify all inputs in all conditions. A safety-critical system must therefore be designed to fail in a manner that preserves safety — for example, by defaulting to a conservative speed or by escalating to human judgement when perception confidence falls below a defined threshold.

The broader lesson from safety-critical systems is that fail-safe behaviour cannot be retrofitted as a policy layer on top of an unsafe architecture. It must be architecturally enforced — embedded in the system's control flow as a structural property, not as a configuration option. A system that can be configured to bypass its own failsafe mechanisms is not fail-safe; it is a system with a failsafe mode that is not always active. The Bailout Protocol adopts this principle: its trigger conditions and escalation paths are defined at the protocol level, with no defined bypass operation that would allow a machine actor to suppress or redirect a bailout trigger.

% -------------------------------------------------------------------------- %
\subsection{The BX3 Framework's Three-Layer Model}
\label{sec:bx3-model}

The BX3 Framework provides the architectural substrate on which the Bailout Protocol is built. The framework defines three immutable functional layers — Purpose Layer, Bounds Engine, and Fact Layer — with enforced properties that hold regardless of which actor occupies each layer. The actor-agnostic, function-centered definition means that any actor satisfying a layer's functional requirements may occupy that layer, including humans, AI systems, and mechanical processes.

The \textbf{Purpose Layer} is the Human Accountability Anchor. It sets strategic goals, defines operational boundaries, and issues authorisation for state transitions. The Purpose Layer is the only layer that can establish a delegation relationship; it is the authoritative origin of any spawned child node's mission and constraints. The requirement that the Purpose Layer remain Human-anchored is non-negotiable and enforced at the architectural level.

The \textbf{Bounds Engine} is the limbless proposer. It computes feasible state transitions, runs simulations, evaluates Bayesian priors, and generates proposals — but it has no physical execution permissions. The Bounds Engine operates in the space of possible actions without the ability to commit any of them. Its outputs are proposals presented to the Fact Layer for deterministic evaluation.

The \textbf{Fact Layer} is the physical firewall. It receives proposed state transitions from the Bounds Engine and executes only those that satisfy pre-defined safety, regulatory, and physical constraints. The Fact Layer is 100\% deterministic: given the same inputs, it must produce identical outputs. It is the only layer that can produce non-reversible physical effects, and its defining characteristic is hard physical constraint enforcement — not software permissions that can be circumvented by privilege escalation.

The three-layer separation is not merely conceptual. Each layer is isolated into a discrete plane; the Bounds Engine never shares a plane with the Fact Layer, and the Fact Layer never shares a plane with the Purpose Layer. This prevents the logic collision failure mode in which a reasoning engine directly executes physical actions without an intermediary constraint evaluation. Loop Isolation — the architectural separation of these three functions — makes logic collision structurally impossible because the constraint evaluation is always performed by a different actor occupying a different layer.

The Bailout Protocol is defined within this architecture as follows: when a child node's Fact Layer encounters a condition that violates its hard constraints and cannot be resolved within its provisioned scope, it generates a trigger event. This trigger propagates up the delegation chain — defined by immutable parent pointers established at spawn time — until it reaches the first Human Accountability Anchor. The critical property is that the escalation chain excludes all machine actors: a bailout trigger cannot be resolved by any intermediate machine actor regardless of its position in the hierarchy. The trigger either reaches a human who can issue a directive, or it exhausts the chain and enters a terminal failure state with a complete forensic ledger of all events that led to the failure.

This architecture ensures that fail-safe behaviour is not a policy configuration but a structural property of the system. The three-layer model provides the clean separation of concerns required for the Bailout Protocol to operate correctly: the Purpose Layer sets the scope of authority, the Bounds Engine operates within that scope and detects boundary violations, the Fact Layer enforces constraints and generates trigger events, and the escalation path terminates only at a human with the authority to issue directives that override the child's current operational context.